% !TEX program = xelatex
% Build: xelatex main.tex, 3 runs (table of contents and cross-references).
% Auxiliary files are written next to the sources, so build a copy elsewhere (build.py does) and keep this folder clean.
% Fonts: Latin Modern for Western text, Fandol (bundled with TeX Live) for Chinese; nothing to install.

% ======== Language of the explainer: \zhtrue for Chinese, \zhfalse for English ========
\newif\ifzh
\zhtrue
% ======================================================================================
% Chinese uses ctexart's Chinese typesetting; English uses the same class with scheme=plain,
% i.e. standard LaTeX typesetting that can still show Chinese characters if needed.
\ifzh
  \documentclass[11pt,a4paper,fontset=fandol]{ctexart}
\else
  \documentclass[11pt,a4paper,fontset=fandol,scheme=plain]{ctexart}
\fi
% \zhen{Chinese}{English} picks the text for the explainer's language. Every fixed label of the
% template goes through it, so switching \zhtrue/\zhfalse above is all a new explainer needs.
\newcommand{\zhen}[2]{\ifzh#1\else#2\fi}
% Placeholder: shown in red until replaced; check_tex.py reports every one left.
\newcommand{\FILL}[1]{\texorpdfstring{\textcolor{red}{[FILL: #1]}}{[FILL: #1]}}

% ======== Per-paper settings ========
\newcommand{\PaperTitle}{\FILL{the paper's title, as published}}
\newcommand{\ExplainerTitle}{\FILL{the title in the explainer's language: a Chinese translation, or the paper's title for English}}
\newcommand{\PaperAuthorsShort}{\FILL{authors' surnames, comma-separated}}
\newcommand{\PaperVenue}{\FILL{venue and year, e.g. NeurIPS 2021}}
\newcommand{\PaperVersion}{\FILL{the version explained, e.g. arXiv v7}}
\newcommand{\PaperHeader}{\FILL{running header, e.g. HaoChen et al.\ (2021)}}
\newcommand{\MadeDate}{\FILL{date YYYY-MM-DD}}
% ====================================

\usepackage[top=2.6cm,bottom=2.4cm,left=2.7cm,right=2.7cm,headheight=15pt]{geometry}
\usepackage{amsmath,amssymb,mathtools,bm}
\usepackage{xcolor}
\usepackage{booktabs,array,tabularx,multirow,longtable}
\usepackage{enumitem}
\usepackage[most]{tcolorbox}
\usepackage{graphicx}
\graphicspath{{figures/}}
\usepackage{tikz}
\usepackage{pgfplots}
\pgfplotsset{compat=1.18}
\usepgfplotslibrary{groupplots}
\usetikzlibrary{arrows.meta,positioning,calc,fit,backgrounds,decorations.pathreplacing,shapes.geometric}
\usepackage{fancyhdr}
\usepackage{needspace}
\usepackage{caption}
\captionsetup{font=small,labelfont={sf,bf}}

% ---------- Colours ----------
\definecolor{headfg}{HTML}{17365D}
\definecolor{linkblue}{HTML}{1F4E9E}
\definecolor{thmfg}{HTML}{1F4E79}  \definecolor{thmbg}{HTML}{EEF4FA}
\definecolor{deffg}{HTML}{0B6E6E}  \definecolor{defbg}{HTML}{EAF6F5}
\definecolor{notefg}{HTML}{9A5B00} \definecolor{notebg}{HTML}{FFF6E5}
\definecolor{chkfg}{HTML}{2E6B2E}  \definecolor{chkbg}{HTML}{EEF7EC}
\definecolor{sumfg}{HTML}{3F3F3F}  \definecolor{sumbg}{HTML}{F2F2F0}
% Chart colours (checked for colour-blind separability; use them in order, don't skip any)
\definecolor{ser1}{HTML}{2A78D6} % blue
\definecolor{ser2}{HTML}{EB6834} % orange
\definecolor{ser3}{HTML}{1BAF7A} % teal green
\definecolor{ser4}{HTML}{EDA100} % yellow
\definecolor{ser5}{HTML}{E87BA4} % magenta
\definecolor{ser6}{HTML}{008300} % green
\definecolor{inkmain}{HTML}{0B0B0B}
\definecolor{inksec}{HTML}{52514E}
\definecolor{inkgrid}{HTML}{E1E0D9}
\definecolor{inkaxis}{HTML}{C3C2B7}

\usepackage[colorlinks=true,linkcolor=linkblue,urlcolor=linkblue,citecolor=linkblue,
  bookmarksnumbered=true,bookmarksopen=true,bookmarksopenlevel=1]{hyperref}
\hypersetup{
  pdftitle={\ExplainerTitle\zhen{——论文详解}{: an explainer}},
  pdfauthor={\zhen{讲解整理：Claude；原论文作者 }{Explainer by Claude; paper by }\PaperAuthorsShort},
  pdfsubject={\PaperTitle\space(\PaperVenue)},
}

% ---------- Headings and layout ----------
\ctexset{
  section={format=\Large\sffamily\bfseries\color{headfg}\raggedright,
           beforeskip=20pt,afterskip=10pt},
  subsection={format=\large\sffamily\bfseries\color{headfg}\raggedright,
           beforeskip=14pt,afterskip=6pt},
  subsubsection={format=\normalsize\sffamily\bfseries\color{headfg}\raggedright,
           beforeskip=10pt,afterskip=4pt},
}
% The explainer's own figures are numbered (插图 N / Diagram N); the paper's figures are captioned
% 原文图 N / Paper Figure N with \caption* and take no number.
\ctexset{figurename={\zhen{插图}{Diagram}}}
\ctexset{part={format=\LARGE\sffamily\bfseries\color{headfg}\centering,beforeskip=22pt,afterskip=12pt}}
\setcounter{secnumdepth}{3}
\setcounter{tocdepth}{1}
\setlist{itemsep=2pt,topsep=3pt,parsep=0pt,leftmargin=2em}
\numberwithin{equation}{section}
\allowdisplaybreaks[3]
\setlength{\emergencystretch}{2em}

\pagestyle{fancy}
\fancyhf{}
\fancyhead[L]{\small\color{inksec}\PaperHeader\ \zhen{论文详解}{explained}}
\fancyhead[R]{\small\color{inksec}\nouppercase{\leftmark}}
\fancyfoot[C]{\small\color{inksec}\thepage}
\renewcommand{\headrulewidth}{0.3pt}
\renewcommand{\sectionmark}[1]{\markboth{#1}{}}
% ctex's \partmark empties the header; show the part's title instead until the next section
\renewcommand{\partmark}[1]{\markboth{#1}{}}

% ---------- Five content boxes: the colour shows where the content comes from ----------
\tcbset{
  boxbase/.style={enhanced jigsaw,breakable,frame hidden,sharp corners,
    left=9pt,right=8pt,top=6pt,bottom=6pt,before skip=9pt,after skip=9pt,
    parbox=false},
}
\newcommand{\boxhead}[2]{\noindent\textcolor{#1}{\sffamily\bfseries #2}\par\smallskip}
% The paper's content: theorems, lemmas, propositions, corollaries, claims, assumptions
\newtcolorbox{thmbox}[1]{boxbase,colback=thmbg,borderline west={2.6pt}{0pt}{thmfg},
  before upper={\boxhead{thmfg}{#1}}}
% The paper's content: definitions
\newtcolorbox{defbox}[1]{boxbase,colback=defbg,borderline west={2.6pt}{0pt}{deffg},
  before upper={\boxhead{deffg}{#1}}}
% The explainer's addition: explanations, connections and assessments (not in the paper)
\newtcolorbox{notebox}[1][\zhen{讲解注}{Explainer's note}]{boxbase,colback=notebg,borderline west={2.6pt}{0pt}{notefg},
  before upper={\boxhead{notefg}{#1}}}
% The explainer's addition: examples and numbers I computed (not in the paper; every number checked by a script)
\newtcolorbox{checkbox}[1][\zhen{验算}{Check}]{boxbase,colback=chkbg,borderline west={2.6pt}{0pt}{chkfg},
  before upper={\boxhead{chkfg}{#1}}}
% Reading guide; only a very long section ends with one listing its points
\newtcolorbox{sumbox}[1][\zhen{小结}{Summary}]{boxbase,colback=sumbg,borderline west={2.6pt}{0pt}{sumfg},
  before upper={\boxhead{sumfg}{#1}}}
\newcommand{\proofhead}[1]{\par\medskip\noindent{\sffamily\bfseries\color{headfg}#1}\quad}
% One plain-words sentence right after a theorem box
\newcommand{\plainsay}[1]{\par\smallskip\noindent{\sffamily\bfseries\color{headfg}\zhen{用大白话说：}{In plain words: }}#1\par\smallskip}
\newcommand{\qedsym}{\hfill$\square$}

% ---------- General notation ----------
% A paper equation number: \ptag{6} shows "(原文 6)" / "(paper 6)" and creates the link target eq:o6
\newcommand{\ptag}[1]{\tag{\zhen{原文}{paper} #1}\label{eq:o#1}}
% Citing a paper equation: 式~\peq{6} / Eq.~\peq{6} shows "式 (6)" / "Eq. (6)", and (6) links to where the explainer displays it
\newcommand{\peq}[1]{\hyperref[eq:o#1]{(#1)}}
\newcommand{\E}{\mathbb{E}}
\newcommand{\R}{\mathbb{R}}
\newcommand{\cX}{\mathcal{X}}
\newcommand{\cF}{\mathcal{F}}
\newcommand{\norm}[1]{\lVert #1\rVert}
\newcommand{\abs}[1]{\lvert #1\rvert}
\newcommand{\ip}[2]{\langle #1,#2\rangle}
\newcommand{\tp}{^{\top}}
\newcommand{\ind}{\mathbf{1}}
\newcommand{\Otil}{\widetilde{O}}
\newcommand{\poly}{\mathrm{poly}}
\DeclareMathOperator*{\argmin}{arg\,min}
\DeclareMathOperator*{\argmax}{arg\,max}
\DeclareMathOperator{\Tr}{Tr}
\DeclareMathOperator{\diag}{diag}

% ---------- This paper's own notation ----------
% e.g. \newcommand{\Loss}{\mathcal{L}}

% ---------- Chart styles ----------
\pgfplotsset{
  cleanplot/.style={
    axis line style={inkaxis,line width=0.6pt},
    tick style={inkaxis,line width=0.4pt},
    grid=major,major grid style={inkgrid,line width=0.4pt},
    every tick label/.append style={font=\footnotesize,color=inksec},
    label style={font=\small,color=inksec},
    title style={font=\small,color=inkmain},
    legend style={font=\footnotesize,draw=none,fill=white,fill opacity=0.9,text opacity=1,
      row sep=-1pt,inner sep=2pt},
    legend cell align=left,
    tick align=outside,
    axis x line*=bottom,axis y line*=left,
  },
  serA/.style={ser1,line width=1.2pt,mark=*,mark size=1.9pt,mark options={solid,draw=white,line width=0.4pt,fill=ser1}},
  serB/.style={ser2,line width=1.2pt,mark=square*,mark size=1.7pt,mark options={solid,draw=white,line width=0.4pt,fill=ser2}},
  serC/.style={ser3,line width=1.2pt,mark=triangle*,mark size=2.4pt,mark options={solid,draw=white,line width=0.4pt,fill=ser3}},
  serD/.style={ser4,line width=1.2pt,mark=diamond*,mark size=2.5pt,mark options={solid,draw=white,line width=0.4pt,fill=ser4}},
  serE/.style={ser5,line width=1.2pt,mark=pentagon*,mark size=2.3pt,mark options={solid,draw=white,line width=0.4pt,fill=ser5}},
  serF/.style={ser6,line width=1.2pt,mark=square*,mark size=1.7pt,mark options={solid,rotate=45,draw=white,line width=0.4pt,fill=ser6}},
}

\begin{document}

% Cover, info table, reading guide, table of contents. Fixed text switches with \zhtrue/\zhfalse in main.tex.
\thispagestyle{plain}
\begin{center}
  {\LARGE\sffamily\bfseries\color{headfg} \ExplainerTitle}\\[6pt]
  {\Large\sffamily\color{headfg} \zhen{论文详解}{An explainer}}\\[8pt]
  \ifzh{\large \PaperTitle}\\[5pt]\fi
  {\normalsize \PaperAuthorsShort\quad \PaperVenue}\\[3pt]
  {\small\color{inksec} \zhen{中文讲解\quad 整理日期}{Prepared} \MadeDate}
\end{center}
\vspace{2pt}

\begin{center}
\small
\begin{tabularx}{\linewidth}{@{}>{\sffamily\bfseries}p{\zhen{4.6em}{5.4em}}>{\raggedright\arraybackslash}X@{}}
\toprule
\zhen{题目}{Title} & \PaperTitle \\
\zhen{作者}{Authors} & \FILL{every author with their affiliation in parentheses, e.g. Jeff Z.\ HaoChen (Stanford University)}\\
\zhen{发表}{Venue} & \FILL{full name and edition of the conference or journal, location or virtual, proceedings volume and pages; oral or spotlight only if verified on the venue's site}\\
\zhen{预印本}{Preprint} & \FILL{arXiv ID; v1 submission date; latest version and its date}\\
\zhen{依据}{Basis} & \FILL{the version explained, and the pages of its main text, references and appendices}\zhen{；与正式版的差异见~\ref{app:versions}}{; differences from the formal version in~\ref{app:versions}}\\
\bottomrule
\end{tabularx}
\end{center}

\begin{sumbox}[\zhen{阅读说明}{How to read this}]
\zhen{\textbf{读者。}假设你有计算机本科的基础；本科课程不教的工具，在第一次用到时讲清。读完这份讲解，不必再打开原文。}{\textbf{Reader.} You are assumed to have a CS undergraduate background; tools that undergraduate courses don't teach are explained where they are first used. After reading this explainer you don't need to open the paper.}

\smallskip
\zhen{\textbf{两种读法。}只读第一部分：第~\ref{sec:plain} 节不用公式讲完整篇论文，后面各节给出全部结果的准确陈述、含义和实验，第~\ref{sec:overview} 节是速查表。要核对每一步推导，再读第二部分的完整证明。}{\textbf{Two ways to read.} Read only Part~I: Section~\ref{sec:plain} tells the whole paper without formulas, the following sections give the exact statement, meaning and experiments of every result, and Section~\ref{sec:overview} is a lookup table. To check every step of the derivations, also read the complete proofs in Part~II.}

\smallskip
\zhen{\textbf{框的含义。}\textcolor{thmfg}{\textbf{蓝框}}（定理、引理、假设）和\textcolor{deffg}{\textbf{青框}}（定义）是原文的陈述，编号同原文 \PaperVersion，只改正了~\ref{app:errata} 列出的笔误。\textcolor{notefg}{\textbf{橙框}}是我补充的解释和评价，\textcolor{chkfg}{\textbf{绿框}}是我算的例子和验算，数字都由脚本核对。无框的正文是我按思路重新组织的论文内容，不是逐句翻译。}{\textbf{What the boxes mean.} \textcolor{thmfg}{\textbf{Blue boxes}} (theorems, lemmas, assumptions) and \textcolor{deffg}{\textbf{teal boxes}} (definitions) are the paper's statements, numbered as in \PaperVersion, with only the typos listed in~\ref{app:errata} corrected. \textcolor{notefg}{\textbf{Orange boxes}} are my explanations and assessments, and \textcolor{chkfg}{\textbf{green boxes}} are examples and checks I computed, every number checked by a script. Unboxed text is the paper's content, reorganized by me along the argument rather than retold sentence by sentence.}

\smallskip
\zhen{\textbf{编号。}公式右边标“(原文 6)”的，就是原文编号为 6 的公式；正文里蓝色的式号可以点击，跳到讲解里展示这个公式的地方。“原文图 N”是原图，“原文表 N”是重排过的原表；有编号的插图是本讲解自己画的。}{\textbf{Numbering.} A formula tagged ``(paper 6)'' on the right is the one numbered (6) in the paper; blue equation numbers in the text are links to where this explainer displays that formula. ``Paper Figure N'' is the paper's original figure and ``Paper Table N'' its table, retyped; numbered diagrams are this explainer's own.}
\end{sumbox}

% Tighter table-of-contents entries; no page break after it, Part I starts right away
\begingroup
\let\origaddvspace\addvspace
\renewcommand{\addvspace}[1]{\origaddvspace{0.5em}}
\tableofcontents
\endgroup


\part{\zhen{轻松读懂这篇论文}{Understanding the paper}}
\section{\zhen{不用公式，先把整篇论文讲一遍}{The whole paper without formulas}}\label{sec:plain}

% One to two pages, with no display formulas and no derivations. One idea per paragraph, usually one paragraph each; undergraduate knowledge needs no explanation.

\paragraph{\zhen{问题。}{Problem.}}\FILL{what problem, told through a concrete scene}

\paragraph{\zhen{已有做法的缺口。}{Gap in existing work.}}\FILL{where existing methods or theory fall short}

\paragraph{\zhen{关键想法。}{Key idea.}}\FILL{the idea, in images and metaphors}

\paragraph{\zhen{主要结论。}{Main results.}}\FILL{the results in words; name one or two key quantities at most}

\paragraph{\zhen{实验。}{Experiments.}}\FILL{what was seen, with the one or two most important numbers}

\paragraph{\zhen{意义与局限。}{Significance and limitations.}}\FILL{why it matters and where it stops}

% If the paper relates to explainers already in the folder, add a paragraph titled \zhen{与已有讲解的关系。}{Relation to earlier explainers.} with one or two sentences per explainer.

% Add in order: 02_prelim (omit if not needed), 03_setup, the main results, experiments, discussion (related work, conclusions and assessment)
\section{\zhen{全部结果速查}{All results at a glance}}\label{sec:overview}

\zhen{这一节只供查找，不含新内容。}{This section is for looking things up; it contains nothing new.}

% Experiment-driven papers: replace "All main results" with a method-components table
% (component | role | formula or algorithm | where in the explainer) and delete "Relations between results".
\subsection{\zhen{所有主要结果}{All main results}}

% The statement column uses the corrected form (as in E2 / 附 2). For a conclusion that is the explainer's own,
% write "(explainer's note)" / "（讲解注）" in the first column.
\begingroup
\small
\renewcommand{\arraystretch}{1.35}
\begin{longtable}{@{}>{\raggedright\arraybackslash}p{0.13\linewidth}>{\raggedright\arraybackslash}p{0.5\linewidth}>{\raggedright\arraybackslash}p{\dimexpr0.37\linewidth-4\tabcolsep\relax}@{}}
\toprule
\textbf{\zhen{结果}{Result}} & \textbf{\zhen{内容}{Statement}} & \textbf{\zhen{含义}{Meaning}} \\
\midrule
\endfirsthead
\toprule
\textbf{\zhen{结果}{Result}} & \textbf{\zhen{内容}{Statement}} & \textbf{\zhen{含义（续）}{Meaning (cont.)}} \\
\midrule
\endhead
\bottomrule
\endlastfoot
\FILL{one row per result} & & \\
\end{longtable}
\endgroup

\subsection{\zhen{结果之间的关系}{Relations between results}}

% Keep only when there are many results (6 or more) with non-trivial dependencies; otherwise delete this subsection.
% TikZ dependency graph: arrows from the lemma used to the theorem that uses it; Part I and Part II results separated by a dashed line.
\FILL{dependency graph}

\subsection{\zhen{符号速查}{Notation}}

\begingroup
\small
\renewcommand{\arraystretch}{1.3}
\begin{longtable}{@{}>{\raggedright\arraybackslash}p{0.3\linewidth}>{\raggedright\arraybackslash}p{\dimexpr0.7\linewidth-2\tabcolsep\relax}@{}}
\toprule
\textbf{\zhen{符号}{Symbol}} & \textbf{\zhen{含义（首次出现）}{Meaning (first appearance)}} \\
\midrule
\endfirsthead
\toprule
\textbf{\zhen{符号}{Symbol}} & \textbf{\zhen{含义（首次出现）（续）}{Meaning (first appearance, cont.)}} \\
\midrule
\endhead
\bottomrule
\endlastfoot
\FILL{one row per symbol} & \\
\end{longtable}
\endgroup


% Parts don't force page breaks, to avoid big blank areas.
% Theory papers keep the "complete proofs" title and intro below; experiment-driven papers use the "full details" pair instead.
\part{\zhen{完整证明（第一次读可以跳过）}{Complete proofs (skip on first reading)}}
\zhen{全部证明都在这里，补全了原文省略的非平凡步骤。原文的笔误直接改正，改了什么见~\ref{app:errata}。}{Every proof is here, with the non-trivial steps the paper omits filled in. Typos in the paper are corrected in place; \ref{app:errata} lists what was changed.}
% \part{\zhen{完整细节（第一次读可以跳过）}{Full details (skip on first reading)}}
% \zhen{全部推导和实验细节都在这里，补全了原文省略的步骤。原文的笔误直接改正，改了什么见~\ref{app:errata}。}{Every derivation and experimental detail is here, with the steps the paper omits filled in. Typos in the paper are corrected in place; \ref{app:errata} lists what was changed.}
% Then add one section per paper appendix, in order.

\part{\zhen{附录}{Appendices}}
% The explainer's own appendices, numbered E1, E2, E3 (Chinese: 附 1, 附 2, 附 3) so they aren't
% confused with the paper's Appendices A, B, C.
\appendix
\ctexset{section/number={\zhen{附}{E}\arabic{section}}}

\section{\FILL{a title naming the versions compared, e.g. \zhen{arXiv v7 与 NeurIPS 正式版的差异}{Differences between arXiv v7 and the NeurIPS version}}}\label{app:versions}

\FILL{one sentence naming the two versions compared: the latest preprint version with its date, and the formal version with its venue and where it and its supplement came from (older preprint versions are not compared). The table lists only substantive differences; delete it if there is no formal version or the two differ only in layout and wording.}

\begingroup
\small
\renewcommand{\arraystretch}{1.3}
\begin{longtable}{@{}>{\raggedright\arraybackslash}p{0.13\linewidth}>{\raggedright\arraybackslash}p{0.43\linewidth}>{\raggedright\arraybackslash}p{\dimexpr0.44\linewidth-4\tabcolsep\relax}@{}}
\toprule
\textbf{\zhen{项目}{Item}} & \textbf{\FILL{latest preprint, e.g. arXiv v7}} & \textbf{\FILL{formal version, e.g. NeurIPS 2021}} \\
\midrule
\endfirsthead
\toprule
\textbf{\zhen{项目}{Item}} & \textbf{\FILL{latest preprint, e.g. arXiv v7}} & \textbf{\FILL{formal version, e.g. NeurIPS 2021}\zhen{（续）}{ (cont.)}} \\
\midrule
\endhead
\bottomrule
\endlastfoot
\zhen{篇幅}{Length} & \FILL{pages} & \FILL{pages} \\
\end{longtable}
\endgroup

\section{\zhen{原文的错误与一致性检查}{Errors and consistency checks in the paper}}\label{app:errata}

\FILL{two or three sentences: whether these problems shake the main conclusions; the category of each entry (conditions or constants, missing steps, substantive gaps, wording and notation); that they were independently reviewed}

% Each entry gives the problem, the correction and how it was verified, and ends with the scope of the impact;
% compact, no reasoning. Ordered by severity.
\begingroup
\small
\renewcommand{\arraystretch}{1.3}
\begin{longtable}{@{}>{\raggedright\arraybackslash}p{0.04\linewidth}>{\raggedright\arraybackslash}p{0.17\linewidth}>{\raggedright\arraybackslash}p{\dimexpr0.79\linewidth-4\tabcolsep\relax}@{}}
\toprule
 & \textbf{\zhen{位置}{Location}} & \textbf{\zhen{问题}{Problem}} \\
\midrule
\endfirsthead
\toprule
 & \textbf{\zhen{位置}{Location}} & \textbf{\zhen{问题（续）}{Problem (cont.)}} \\
\midrule
\endhead
\bottomrule
\endlastfoot
1 & \FILL{location} & \FILL{problem; correction; how it was verified; scope} \\
\end{longtable}
\endgroup

\paragraph{\zhen{核对过没有问题的部分。}{Parts checked and found correct.}}\FILL{the main results verified by script with no problem found}

\section{\zhen{原文参考文献及其作用}{The paper's references and their roles}}\label{app:refs}

\FILL{all N references of the version explained. Those whose results, methods, data or code the paper uses get one row each, in the paper's order; those only listed in batches in the introduction or related work are grouped by role at the end. Mark those new relative to the formal version.}

\begingroup
\small
\renewcommand{\arraystretch}{1.22}
\begin{longtable}{@{}>{\raggedright\arraybackslash}p{27.5em}>{\raggedright\arraybackslash}p{13.5em}@{}}
\toprule
\textbf{\zhen{文献}{Reference}} & \textbf{\zhen{在文中的作用}{Role in the paper}} \\
\midrule
\endfirsthead
\toprule
\textbf{\zhen{文献}{Reference}} & \textbf{\zhen{在文中的作用（续）}{Role in the paper (cont.)}} \\
\midrule
\endhead
\bottomrule
\endlastfoot
\FILL{reference} & \FILL{role} \\
% References cited only in batches go below this rule, one row per group; delete the rule and heading if there are none.
\midrule
\multicolumn{2}{@{}l@{}}{\textbf{\zhen{只在引言和相关工作里成批列举的文献，按作用分组}{Cited only in batches in the introduction or related work, grouped by role}}} \\
\FILL{the group's references: authors, year, and what each is} & \FILL{the group's role} \\
\end{longtable}
\endgroup


\end{document}
